\documentclass[12pt]{ximera}
\typeout{Start loading xmPreamble.tex}%

% Add here extra macro's that are loaded automatically by all documents of claas 'ximera' or 'xourse' in this repo

%%
%%  Example:
%%
% \newcommand{\R}{\mathbb{R}}
\usepackage{tikz}
\usetikzlibrary{positioning, arrows.meta}
\usepackage{multicol}
\usepackage{enumitem}
\usepackage{caption}
\usepackage{amsmath}
\usepackage{amsthm}
\usepackage{fontenc}

%% ---------------------------------------------------------------
%% Shared worksheet macros.
%%
%% These came from the Summer 2026 worksheet set, where each worksheet carried its
%% own copy. They have to live here instead: a xourse inlines every activity into a
%% single document, so duplicate \newcommand definitions across activities collide,
%% and the xourse gobbles \input so a per-topic macro file never loads.
%%
%% They use \providecommand, NOT \newcommand. ximera.cls loads this file automatically,
%% and every activity also carries an explicit \typeout{Start loading xmPreamble.tex}%

% Add here extra macro's that are loaded automatically by all documents of claas 'ximera' or 'xourse' in this repo

%%
%%  Example:
%%
% \newcommand{\R}{\mathbb{R}}
\usepackage{tikz}
\usetikzlibrary{positioning, arrows.meta}
\usepackage{multicol}
\usepackage{enumitem}
\usepackage{caption}
\usepackage{amsmath}
\usepackage{amsthm}
\usepackage{fontenc}

%% ---------------------------------------------------------------
%% Shared worksheet macros.
%%
%% These came from the Summer 2026 worksheet set, where each worksheet carried its
%% own copy. They have to live here instead: a xourse inlines every activity into a
%% single document, so duplicate \newcommand definitions across activities collide,
%% and the xourse gobbles \input so a per-topic macro file never loads.
%%
%% They use \providecommand, NOT \newcommand. ximera.cls loads this file automatically,
%% and every activity also carries an explicit \typeout{Start loading xmPreamble.tex}%

% Add here extra macro's that are loaded automatically by all documents of claas 'ximera' or 'xourse' in this repo

%%
%%  Example:
%%
% \newcommand{\R}{\mathbb{R}}
\usepackage{tikz}
\usetikzlibrary{positioning, arrows.meta}
\usepackage{multicol}
\usepackage{enumitem}
\usepackage{caption}
\usepackage{amsmath}
\usepackage{amsthm}
\usepackage{fontenc}

%% ---------------------------------------------------------------
%% Shared worksheet macros.
%%
%% These came from the Summer 2026 worksheet set, where each worksheet carried its
%% own copy. They have to live here instead: a xourse inlines every activity into a
%% single document, so duplicate \newcommand definitions across activities collide,
%% and the xourse gobbles \input so a per-topic macro file never loads.
%%
%% They use \providecommand, NOT \newcommand. ximera.cls loads this file automatically,
%% and every activity also carries an explicit \input{../xmPreamble}, so this file is
%% read twice. That was harmless while it held only \usepackage lines (LaTeX skips a
%% package it has already loaded) but a second \newcommand is a hard error.
%%
%%   \blank[width]        fill-in-the-blank rule (default 2.5cm)
%%   \showwork            "show and explain your work" reminder
%%   \showworkline        ... plus which schedule line was used
%%   \showworkyear        ... plus which year's schedule was used
%%   \schedhead{...}      centred bold caption above a tiered-rate table
%%   \samesquares[scale]{left}{right}   two equal-area squares, labelled
%%   \samecubes[scale]{left}{right}     two equal-volume cubes, labelled
%% ---------------------------------------------------------------

\providecommand{\blank}[1][2.5cm]{\underline{\hspace{#1}}}
\providecommand{\showwork}{\textbf{REMEMBER TO SHOW AND EXPLAIN YOUR WORK.}}
\providecommand{\showworkline}{\textbf{REMEMBER TO SHOW AND EXPLAIN YOUR WORK.
  Explanation should include how and why you know which line of the
  schedule was used to calculate this amount.}}
\providecommand{\showworkyear}{\begin{sloppypar}\textbf{REMEMBER TO SHOW AND
  EXPLAIN YOUR WORK. Explanation should include how and why you know which
  line of the year's tax schedule was used to calculate this tax
  amount.}\end{sloppypar}}
\providecommand{\schedhead}[1]{\begin{center}\textbf{#1}\end{center}}

\providecommand{\samesquares}[3][1]{%
\begin{center}
{\footnotesize These squares are the \textbf{same size}.}

\vspace{1.5mm}
\begin{tikzpicture}[scale=#1,every node/.style={scale=#1}]
  \draw (0,0) rectangle (2.4,2.4);
  \node[anchor=north west] at (0.15,2.25) {Area:};
  \node[anchor=east]  at (-0.15,1.2) {#2};
  \node[anchor=north] at (1.2,-0.15) {#2};
  \begin{scope}[xshift=4.6cm]
    \draw (0,0) rectangle (2.4,2.4);
    \node[anchor=north west] at (0.15,2.25) {Area:};
    \node[anchor=east]  at (-0.15,1.2) {#3};
    \node[anchor=north] at (1.2,-0.15) {#3};
  \end{scope}
\end{tikzpicture}
\end{center}}

\providecommand{\samecubes}[3][1]{%
\begin{center}
{\footnotesize These cubes are the \textbf{same size}.}

\vspace{1.5mm}
\begin{tikzpicture}[scale=#1,every node/.style={scale=#1}]
  \foreach \dx in {0,5.2} {
    \begin{scope}[xshift=\dx cm]
      \draw (0,0) rectangle (2.0,2.0);
      \draw (0,2.0) -- (0.65,2.6) -- (2.65,2.6) -- (2.0,2.0);
      \fill[black!12] (2.0,0) -- (2.65,0.6) -- (2.65,2.6) -- (2.0,2.0) -- cycle;
      \draw (2.0,0) -- (2.65,0.6) -- (2.65,2.6) -- (2.0,2.0);
      \node[anchor=north west] at (0.1,1.9) {Volume:};
    \end{scope}
  }
  \node[anchor=east]  at (-0.15,1.0) {#2};
  \node[anchor=north] at (1.0,-0.15) {#2};
  \node[anchor=west]  at (2.25,0.4) {#2};
  \node[anchor=east]  at (5.05,1.0) {#3};
  \node[anchor=north] at (6.2,-0.15) {#3};
  \node[anchor=west]  at (7.45,0.4) {#3};
\end{tikzpicture}
\end{center}}

%% NOTE: do NOT add \usepackage to individual activity .tex files.
%% When a xourse inlines an activity it gobbles \documentclass and \input, but a bare
%% \usepackage still executes -- after \begin{document} -- and errors out.
%% All shared packages and macros belong here.
%%
%% ximera.cls already loads: amsmath, amssymb, amsthm, cancel, enumitem, graphicx,
%% hyperref, listings, pgfplots, tikz, url, xcolor. No need to re-load those.
, so this file is
%% read twice. That was harmless while it held only \usepackage lines (LaTeX skips a
%% package it has already loaded) but a second \newcommand is a hard error.
%%
%%   \blank[width]        fill-in-the-blank rule (default 2.5cm)
%%   \showwork            "show and explain your work" reminder
%%   \showworkline        ... plus which schedule line was used
%%   \showworkyear        ... plus which year's schedule was used
%%   \schedhead{...}      centred bold caption above a tiered-rate table
%%   \samesquares[scale]{left}{right}   two equal-area squares, labelled
%%   \samecubes[scale]{left}{right}     two equal-volume cubes, labelled
%% ---------------------------------------------------------------

\providecommand{\blank}[1][2.5cm]{\underline{\hspace{#1}}}
\providecommand{\showwork}{\textbf{REMEMBER TO SHOW AND EXPLAIN YOUR WORK.}}
\providecommand{\showworkline}{\textbf{REMEMBER TO SHOW AND EXPLAIN YOUR WORK.
  Explanation should include how and why you know which line of the
  schedule was used to calculate this amount.}}
\providecommand{\showworkyear}{\begin{sloppypar}\textbf{REMEMBER TO SHOW AND
  EXPLAIN YOUR WORK. Explanation should include how and why you know which
  line of the year's tax schedule was used to calculate this tax
  amount.}\end{sloppypar}}
\providecommand{\schedhead}[1]{\begin{center}\textbf{#1}\end{center}}

\providecommand{\samesquares}[3][1]{%
\begin{center}
{\footnotesize These squares are the \textbf{same size}.}

\vspace{1.5mm}
\begin{tikzpicture}[scale=#1,every node/.style={scale=#1}]
  \draw (0,0) rectangle (2.4,2.4);
  \node[anchor=north west] at (0.15,2.25) {Area:};
  \node[anchor=east]  at (-0.15,1.2) {#2};
  \node[anchor=north] at (1.2,-0.15) {#2};
  \begin{scope}[xshift=4.6cm]
    \draw (0,0) rectangle (2.4,2.4);
    \node[anchor=north west] at (0.15,2.25) {Area:};
    \node[anchor=east]  at (-0.15,1.2) {#3};
    \node[anchor=north] at (1.2,-0.15) {#3};
  \end{scope}
\end{tikzpicture}
\end{center}}

\providecommand{\samecubes}[3][1]{%
\begin{center}
{\footnotesize These cubes are the \textbf{same size}.}

\vspace{1.5mm}
\begin{tikzpicture}[scale=#1,every node/.style={scale=#1}]
  \foreach \dx in {0,5.2} {
    \begin{scope}[xshift=\dx cm]
      \draw (0,0) rectangle (2.0,2.0);
      \draw (0,2.0) -- (0.65,2.6) -- (2.65,2.6) -- (2.0,2.0);
      \fill[black!12] (2.0,0) -- (2.65,0.6) -- (2.65,2.6) -- (2.0,2.0) -- cycle;
      \draw (2.0,0) -- (2.65,0.6) -- (2.65,2.6) -- (2.0,2.0);
      \node[anchor=north west] at (0.1,1.9) {Volume:};
    \end{scope}
  }
  \node[anchor=east]  at (-0.15,1.0) {#2};
  \node[anchor=north] at (1.0,-0.15) {#2};
  \node[anchor=west]  at (2.25,0.4) {#2};
  \node[anchor=east]  at (5.05,1.0) {#3};
  \node[anchor=north] at (6.2,-0.15) {#3};
  \node[anchor=west]  at (7.45,0.4) {#3};
\end{tikzpicture}
\end{center}}

%% NOTE: do NOT add \usepackage to individual activity .tex files.
%% When a xourse inlines an activity it gobbles \documentclass and \input, but a bare
%% \usepackage still executes -- after \begin{document} -- and errors out.
%% All shared packages and macros belong here.
%%
%% ximera.cls already loads: amsmath, amssymb, amsthm, cancel, enumitem, graphicx,
%% hyperref, listings, pgfplots, tikz, url, xcolor. No need to re-load those.
, so this file is
%% read twice. That was harmless while it held only \usepackage lines (LaTeX skips a
%% package it has already loaded) but a second \newcommand is a hard error.
%%
%%   \blank[width]        fill-in-the-blank rule (default 2.5cm)
%%   \showwork            "show and explain your work" reminder
%%   \showworkline        ... plus which schedule line was used
%%   \showworkyear        ... plus which year's schedule was used
%%   \schedhead{...}      centred bold caption above a tiered-rate table
%%   \samesquares[scale]{left}{right}   two equal-area squares, labelled
%%   \samecubes[scale]{left}{right}     two equal-volume cubes, labelled
%% ---------------------------------------------------------------

\providecommand{\blank}[1][2.5cm]{\underline{\hspace{#1}}}
\providecommand{\showwork}{\textbf{REMEMBER TO SHOW AND EXPLAIN YOUR WORK.}}
\providecommand{\showworkline}{\textbf{REMEMBER TO SHOW AND EXPLAIN YOUR WORK.
  Explanation should include how and why you know which line of the
  schedule was used to calculate this amount.}}
\providecommand{\showworkyear}{\begin{sloppypar}\textbf{REMEMBER TO SHOW AND
  EXPLAIN YOUR WORK. Explanation should include how and why you know which
  line of the year's tax schedule was used to calculate this tax
  amount.}\end{sloppypar}}
\providecommand{\schedhead}[1]{\begin{center}\textbf{#1}\end{center}}

\providecommand{\samesquares}[3][1]{%
\begin{center}
{\footnotesize These squares are the \textbf{same size}.}

\vspace{1.5mm}
\begin{tikzpicture}[scale=#1,every node/.style={scale=#1}]
  \draw (0,0) rectangle (2.4,2.4);
  \node[anchor=north west] at (0.15,2.25) {Area:};
  \node[anchor=east]  at (-0.15,1.2) {#2};
  \node[anchor=north] at (1.2,-0.15) {#2};
  \begin{scope}[xshift=4.6cm]
    \draw (0,0) rectangle (2.4,2.4);
    \node[anchor=north west] at (0.15,2.25) {Area:};
    \node[anchor=east]  at (-0.15,1.2) {#3};
    \node[anchor=north] at (1.2,-0.15) {#3};
  \end{scope}
\end{tikzpicture}
\end{center}}

\providecommand{\samecubes}[3][1]{%
\begin{center}
{\footnotesize These cubes are the \textbf{same size}.}

\vspace{1.5mm}
\begin{tikzpicture}[scale=#1,every node/.style={scale=#1}]
  \foreach \dx in {0,5.2} {
    \begin{scope}[xshift=\dx cm]
      \draw (0,0) rectangle (2.0,2.0);
      \draw (0,2.0) -- (0.65,2.6) -- (2.65,2.6) -- (2.0,2.0);
      \fill[black!12] (2.0,0) -- (2.65,0.6) -- (2.65,2.6) -- (2.0,2.0) -- cycle;
      \draw (2.0,0) -- (2.65,0.6) -- (2.65,2.6) -- (2.0,2.0);
      \node[anchor=north west] at (0.1,1.9) {Volume:};
    \end{scope}
  }
  \node[anchor=east]  at (-0.15,1.0) {#2};
  \node[anchor=north] at (1.0,-0.15) {#2};
  \node[anchor=west]  at (2.25,0.4) {#2};
  \node[anchor=east]  at (5.05,1.0) {#3};
  \node[anchor=north] at (6.2,-0.15) {#3};
  \node[anchor=west]  at (7.45,0.4) {#3};
\end{tikzpicture}
\end{center}}

%% NOTE: do NOT add \usepackage to individual activity .tex files.
%% When a xourse inlines an activity it gobbles \documentclass and \input, but a bare
%% \usepackage still executes -- after \begin{document} -- and errors out.
%% All shared packages and macros belong here.
%%
%% ximera.cls already loads: amsmath, amssymb, amsthm, cancel, enumitem, graphicx,
%% hyperref, listings, pgfplots, tikz, url, xcolor. No need to re-load those.

\title{Units: Change Measurement, Not Quantity}
\begin{document}
\begin{abstract}
Unit conversion as multiplication by a disguised 1, why squared and cubed units
convert by the square and cube of the linear factor, scaling as its own kind of
equivalence, and non-standard equivalences such as price and material density.
\end{abstract}
\maketitle

\section*{Change Measurement, Not Quantity}

Is 12 the same as 1? As a numerical value, no, 12 and 1 are not the same. However, if
we give these values units, they can represent the same thing! For instance, 12 inches
is the same length as 1 foot, and 12 months is the same span as 1 year.

Units are often used in everyday life, even if we don't necessarily think of them as
units. For example, you wouldn't say ``I went to the grocery store and bought a
dozen.'' A dozen what? You've given how many you've bought (the value), but not how
many of what you have bought (the units). Instead, you would say ``I went to the
grocery store and bought a dozen eggs.''

Another aspect to note is that when we convert between units, we are \textbf{not
changing the quantity}, we are just changing \textbf{how we are measuring it}. For
example, if you have 1 gallon or 4 quarts of some liquid, you have the same amount of
liquid since 1 gallon $=$ 4 quarts. The difference is only in how you are choosing to
measure the liquid: gallons (a larger measurement) or quarts (a smaller measurement).
This change of measurement, not quantity, is why we can say that 3 gallons is equal to
$3 \cdot 4 = 12$ quarts. Why do we multiply by 4? Why not divide by 4? We can use the
math of fractions to help us ``see'' why this works.

Since 1 gallon $=$ 4 quarts, multiplying by $\dfrac{4\text{ quarts}}{1\text{ gallon}}$
or by $\dfrac{1\text{ gallon}}{4\text{ quarts}}$ is the same as multiplying the
quantity by 1, since the quantities 4 quarts and 1 gallon are equal, and equal things
divided by each other equal 1:
\[
1 \;=\; \frac{3}{3} \;=\; \frac{138}{138} \;=\; \frac{\pi}{\pi}
\;=\; \frac{4\text{ quarts}}{1\text{ gallon}} \;=\; \frac{1\text{ gallon}}{4\text{ quarts}}
\]
However, we are changing \emph{how} we are measuring it, which is why the numerical
value is changing. Remember, we are only changing the units, not the amount. Hence:
\[
3\text{ gallons} \cdot \frac{4\text{ quarts}}{1\text{ gallon}} \;=\; 12\text{ quarts}
\]

Just as with ``normal'' fractions: when we have the same value in the numerator and the
denominator, they ``cancel away.'' In a similar way of thinking, we can see that there
are ``gallons'' in the numerator of our original value, and ``gallons'' in the
denominator of what we're multiplying by, and thus the ``gallon'' units ``cancel'' with
each other, leaving the ``quart'' units in our answer. It is in needing these like
units to cancel that we can know we would NOT do
\[
3\text{ gallons} \cdot \frac{1\text{ gallon}}{4\text{ quarts}} \;=\; 0.75\text{ square gallons per quart}
\]
As this gives the units ``square gallons per quart'' which is nonsensical, or at least
not something helpful to us for understanding in the real world.

\section*{Units \& Dimensions}

First, let's practice some common conversions using this ``fractional'' approach
described above.

\begin{problem}
How many yards are in a mile?

\hspace{1em} 3 feet $=$ 1 yard \qquad 5,280 feet $=$ 1 mile

There are $\answer{1760}$ yards in a mile.

\begin{freeResponse}
\end{freeResponse}

\begin{solution}
We ``have'' 1 mile, and we are trying to change our measurement of this length into
yards.
\begin{gather*}
1\text{ mile} \cdot \frac{5{,}280\text{ feet}}{1\text{ mile}} \cdot \frac{1\text{ yard}}{3\text{ feet}}
= \frac{5{,}280}{3}\text{ yards} = 1{,}760\text{ yards}
\end{gather*}
Hence there are 1,760 yards in a mile.

\textbf{Note:} It is highly recommended that you write an ``intro'' sentence and a
``concluding'' sentence for each of your answers, as that will help you to provide more
explanation and context to the math of the fractional calculations. This is also a good
tip to keep in mind for other problems where you are not sure how to ``explain more.''
\end{solution}
\end{problem}

\begin{problem}
How many cups of gasoline could fit into a 15 gallon tank?

\hspace{1em} 4 cups $=$ 1 quart \qquad 4 quarts $=$ 1 gallon

A 15 gallon tank holds $\answer{240}$ cups.

\begin{freeResponse}
\end{freeResponse}

\begin{solution}
We have 15 gallons, which we want to convert into cup measurements.
\begin{gather*}
15\text{ gal} \cdot \frac{4\text{ quarts}}{1\text{ gal}} \cdot \frac{4\text{ cups}}{1\text{ quart}}
= 15 \cdot 4 \cdot 4\text{ cups} = 240\text{ cups}
\end{gather*}
So there are 240 cups in 15 gallons.
\end{solution}
\end{problem}

\begin{problem}
The definition of an acre is that 1 acre $=$ 43,560 sq ft. Consider the question:
``100,000 feet is the same as how many acres?'' Explain why this question does not make
sense. In your explanation, show an attempt at answering the question, and explain what
the resulting units would be.

\begin{freeResponse}
\end{freeResponse}

\begin{solution}
This does not make sense because the question is asking us to convert feet, which
measures length, into acres, which measures area, which are not the same thing. Acres
measure the area which SQUARE feet also measure, not ``plain'' feet.

The question might be asked thinking that we could do this:
\begin{gather*}
100{,}000\text{ feet} \cdot \frac{1\text{ acre}}{43{,}560\text{ sq ft}}
\end{gather*}
But this would not ``cancel away'' the sq ft in the denominator, so we would have
\begin{gather*}
100{,}000\text{ feet} \cdot \frac{1\text{ acre}}{43{,}560\text{ sq ft}}
= \frac{100{,}000}{43{,}560}\ \text{acre/feet} = \text{about } 2.3\ \text{acre/feet}.
\end{gather*}
Perhaps acres per foot make sense in some context, but I am not sure of one, and either
way, it would not be converting feet to acres!
\end{solution}
\end{problem}

Just as in the questions above, you will be provided with conversion values that you
will need for problems. However, you will only be given the conversion factors ``12 in
$=$ 1 ft'' or ``5,280 ft $=$ 1 mile'' for a problem which asks about square inches,
square feet, square miles, or cubic inches, cubic feet, or cubic miles. Let's determine
how we should use such conversion factors for the square or cubic versions.

\section*{Determining Dimensional Conversion Equivalences}

In order to convert between units, we need some sort of equivalence. Just like 4 quarts
$=$ 1 gallon or 16 oz $=$ 1 pound are equivalences that we can then use to convert from
quarts to gallons or vice versa, or ounces to pounds or vice versa. Let's determine the
equivalence between square feet and square inches.

\samesquares{1 ft}{\blank[1.1cm]\ in}

\begin{problem}
Calculate the area of the left square using the given ``feet'' units. Then label the
right square with the equivalent quantity in ``inches'' units and calculate its area.
What is the case about these two areas?

1 sq ft $= \answer{144}$ sq inches.

\begin{freeResponse}
\end{freeResponse}

\begin{solution}
Since these squares are the same size, that means they must cover the same area. Since
the left one is 1 ft $\cdot$ 1 ft $=$ 1 sq ft, and the right one is the same, but
measured in inches, and thus an area of 12 inches $\cdot$ 12 inches $=$ 144 sq inches,
this means that 1 sq ft $=$ 144 sq inches!
\end{solution}
\end{problem}

\begin{problem}
Convert 0.75 sq ft into square inches.

$0.75$ sq ft $= \answer{108}$ sq inches.

\begin{solution}
We have 0.75 sq ft and we want to know how many sq inches this area is, so we can use
the fact that 1 sq ft $= 12 \cdot 12 = 144$ square inches to convert it to square
inches.
\begin{gather*}
0.75\text{ sq ft} \cdot \frac{144\text{ sq in}}{1\text{ sq ft}} = 108\text{ sq inches}
\end{gather*}
Hence 0.75 sq ft is the same as 108 sq inches.
\end{solution}
\end{problem}

\begin{problem}
Convert 36,000 square inches into square feet.

$36{,}000$ sq inches $= \answer{250}$ sq feet.

\begin{solution}
We have 36,000 sq inches and we want to know how many sq feet this area is, so we can
use the fact that 1 sq ft $= 12 \cdot 12 = 144$ square inches to convert it to square
feet.
\begin{gather*}
36{,}000\text{ sq inches} \cdot \frac{1\text{ sq ft}}{144\text{ sq inches}} = 250\text{ sq feet}
\end{gather*}
Hence 36,000 sq inches is the same as 250 sq feet.
\end{solution}
\end{problem}

\begin{problem}
Using the idea we established in the previous problems, fill in the following
equivalences.

\hspace{1em} 12 inches $=$ 1 foot \qquad 5,280 feet $=$ 1 mile

\begin{itemize}
\item 1 sq foot $= \answer{144}$ sq inches
\item 1 sq mile $= \answer{27878400}$ sq feet
\item 1 cubic foot $= \answer{1728}$ cubic inches
\item 1 cubic mile $= \answer{147197952000}$ cubic feet
\end{itemize}

\samesquares[0.60]{1 ft}{\blank[1.1cm]\ in}

\samecubes[0.60]{1 ft}{\blank[1.1cm]\ in}

\begin{solution}
\begin{gather*}
1\text{ sq ft} = 1\text{ ft} \cdot 1\text{ ft} = 12\text{ in} \cdot 12\text{ in}
= 12^{2} = 144\text{ sq inches} \\
1\text{ sq mi} = 1\text{ mi} \cdot 1\text{ mi} = 5{,}280\text{ ft} \cdot 5{,}280\text{ ft}
= (5{,}280)^{2} = 27{,}878{,}400\text{ sq ft} \\
1\text{ cu ft} = 1\text{ ft} \cdot 1\text{ ft} \cdot 1\text{ ft}
= 12\text{ in} \cdot 12\text{ in} \cdot 12\text{ in} = 12^{3} = 1{,}728\text{ cu inches} \\
1\text{ cu mi} = 5{,}280\text{ ft} \cdot 5{,}280\text{ ft} \cdot 5{,}280\text{ ft}
= 5{,}280^{3} = 147{,}197{,}952{,}000\text{ cu ft}
\end{gather*}
\end{solution}
\end{problem}

The conversion values you have determined above will \textbf{not be provided to you}.
It is expected that you use provided information such as ``1 foot $=$ 12 inches'' to
determine an equivalence between square inches and square feet, or cubic inches and
cubic feet as needed. Use these equivalences to help you answer the following problems.

\begin{problem}
A lot of land measures 50,000 feet by 150,000 feet. How many square miles is this lot?
Answer this question by first converting the feet into miles.

The lot is about $\answer[tolerance=2]{269.8}$ square miles.

\begin{freeResponse}
\end{freeResponse}

\begin{solution}
If we first convert the feet into miles, we can then multiply the mile lengths to get
the area in sq mile units:
\begin{gather*}
50{,}000\text{ ft} \cdot \frac{1\text{ mile}}{5{,}280\text{ ft}} = \frac{50{,}000}{5{,}280}\text{ miles} = \text{about } 9.5\text{ miles} \\
150{,}000\text{ ft} \cdot \frac{1\text{ mile}}{5{,}280\text{ ft}} = \frac{150{,}000}{5{,}280}\text{ miles} = \text{about } 28.4\text{ miles}
\end{gather*}
Now we multiply the lengths of the land in miles to find the square miles:
\begin{gather*}
9.5\text{ miles} \cdot 28.4\text{ miles} = 269.8\text{ sq miles}
\end{gather*}
The lot of land is about 269.8 sq miles.
\end{solution}
\end{problem}

\begin{problem}
Answer the same question --- how many square miles is the 50,000 ft by 150,000 ft lot
--- but this time by first calculating the square feet of the land.

\begin{freeResponse}
\end{freeResponse}

\begin{solution}
We can calculate the square feet the lot of land is and then convert those sq ft into
sq miles using the fact that 1 sq mi $= 5{,}280^{2} = 27{,}878{,}400$ sq ft.

Area in sq ft:
\begin{gather*}
50{,}000\text{ ft} \cdot 150{,}000\text{ ft} = 7{,}500{,}000{,}000\text{ sq ft}
\end{gather*}
Now convert the area into sq miles:
\begin{gather*}
7{,}500{,}000{,}000\text{ sq ft} \cdot \frac{1\text{ sq mile}}{27{,}878{,}400\text{ sq ft}} = \text{about } 269\text{ sq miles}
\end{gather*}
The lot of land is about 269 sq miles.

The two methods differ slightly (269.8 versus 269) only because the first rounded the
mile lengths to one decimal place before multiplying.
\end{solution}
\end{problem}

\begin{problem}
How many acres is that same lot?

\hspace{1em} 43,560 sq ft $=$ 1 acre

The lot is about $\answer[tolerance=500]{172176}$ acres.

\begin{solution}
Since we calculated how many sq feet the land is, we know it is 7,500,000,000 sq ft. So
we can use the conversion factor to change this into acre measurements:
\begin{gather*}
7{,}500{,}000{,}000\text{ sq ft} \cdot \frac{1\text{ acre}}{43{,}560\text{ sq ft}} = \text{about } 172{,}176\text{ acres}
\end{gather*}
The lot of land is about 172,176 acres.
\end{solution}
\end{problem}

\begin{problem}
A fish tank measures 36 inches by 18 inches by 18 inches. How many cubic feet of water
does it hold? Solve this in two ways, and remember to state what the ``method'' is you
are going to use before each one.

\hspace{1em} 12 inches $=$ 1 foot

The tank holds $\answer[tolerance=0.1]{6.75}$ cubic feet.

\begin{freeResponse}
\end{freeResponse}

\begin{solution}
\textbf{Method 1.} If we first convert the inches into feet, we can then multiply the
feet lengths to get the volume in cubic feet units:
\begin{gather*}
36\text{ inches} \cdot \frac{1\text{ ft}}{12\text{ inches}} = \frac{36}{12}\text{ ft} = 3\text{ ft} \\
18\text{ inches} \cdot \frac{1\text{ ft}}{12\text{ inches}} = \frac{18}{12}\text{ ft} = 1.5\text{ ft}
\end{gather*}
Now we multiply the lengths of the tank in feet to find the cubic feet:
\begin{gather*}
3\text{ ft} \cdot 1.5\text{ ft} \cdot 1.5\text{ ft} = 6.75\text{ cubic feet}
\end{gather*}

\textbf{Method 2.} We can calculate the cubic inches the tank holds and then convert
those cubic inches into cubic feet using the fact that 1 cu ft $= 12^{3} = 1{,}728$ cu
in.
\begin{gather*}
36\text{ in} \cdot 18\text{ in} \cdot 18\text{ in} = 11{,}664\text{ cu in} \\
11{,}664\text{ cu in} \cdot \frac{1\text{ cu ft}}{1{,}728\text{ cu in}} = 6.75\text{ cu ft}
\end{gather*}
Either way, the tank is 6.75 cubic feet.
\end{solution}
\end{problem}

\begin{problem}
How many gallons does the tank hold?

\hspace{1em} 7.48 gallons $=$ 1 cubic foot

The tank holds about $\answer[tolerance=0.5]{50.5}$ gallons.

\begin{solution}
Using either method, we end up with 6.75 cu ft for the tank. So now we just convert
this to gallons using the conversion factor:
\begin{gather*}
6.75\text{ cu ft} \cdot \frac{7.48\text{ gallons}}{1\text{ cu ft}} = 6.75 \cdot 7.48\text{ gallons} = 50.49\text{ gallons}
\end{gather*}
The tank holds about 50.5 gallons.
\end{solution}
\end{problem}

\begin{problem}
Suppose a lake has a surface area of about 748 acres and an average depth of about 18
feet. How many square miles does the lake cover?

\hspace{1em} 43,560 sq ft $=$ 1 acre \qquad 5,280 feet $=$ 1 mile

The lake covers about $\answer[tolerance=0.05]{1.17}$ square miles.

\begin{solution}
The lake covers 748 acres, so we need to convert this to sq ft. Then we can convert the
sq ft to sq miles using the fact that 1 sq mi $= 5{,}280^{2} = 27{,}878{,}400$ sq ft.
\begin{gather*}
748\text{ acres} \cdot \frac{43{,}560\text{ sq ft}}{1\text{ acre}} = 32{,}582{,}880\text{ sq ft} \\
32{,}582{,}880\text{ sq ft} \cdot \frac{1\text{ sq mi}}{27{,}878{,}400\text{ sq ft}}
= \text{about } 1.17\text{ sq miles}
\end{gather*}
The lake covers about 1.17 sq miles of area.
\end{solution}
\end{problem}

\begin{problem}
How many cubic feet are in the lake?

\begin{freeResponse}
\end{freeResponse}

\begin{solution}
Above we found that the lake covers an area of 32,582,880 sq ft. We know the lake has
an average depth of 18 ft. If we multiply the sq ft of area by the ft of depth, we'll
have sq ft $\cdot$ ft $=$ ft $\cdot$ ft $\cdot$ ft $=$ cubic feet, which is what we
want:
\begin{gather*}
32{,}582{,}880\text{ sq ft} \cdot 18\text{ ft} = 586{,}491{,}840\text{ cubic feet}
\end{gather*}
The lake is about 586,491,840 cubic feet.
\end{solution}
\end{problem}

\section*{Units \& Scaling}

When we scale something down or scale something up, we \textbf{are} changing the
quantity. For example, if we scale up by 2 something that is 3 ft, the resulting scaled
up item would be $3\text{ ft} \cdot 2 = 6$ ft. However, if we establish an equivalence
between two measurements to define our scaling, then we can treat it like a ``normal''
unit conversion in order to calculate scaled values.

For example, suppose we have something that is 1 ft, and we want to scale it up to 4
yards. So our equivalence is that 1 ft original $=$ 4 yards scaled. So if we have
something that is 5 ft long that we want to scale up by the same amount, we can use
this equivalence to find out how long the 5 foot item should be when scaled up:
\[
5\text{ ft original} \cdot \frac{4\text{ yards scaled}}{1\text{ ft original}} \;=\; 20\text{ yards scaled}
\]

We could also use this if we were trying to scale something down. Suppose we had 100
yards in our scaled up version. What would be the original length in the un-scaled
version?
\[
100\text{ yards scaled} \cdot \frac{1\text{ ft original}}{4\text{ yards scaled}} \;=\; \text{about } 25\text{ feet original}
\]

Notice the use of the words ``original'' and ``scaled'' in the work above. This is there
because we \textbf{are} changing the quantity and so we need to keep track of whether we
have or want our original quantity or our scaled quantity.

\begin{problem}
Using this same equivalence (1 ft original $=$ 4 yards scaled): suppose we have
something that is 2 yards (original) long. How long should the scaled version of this
be?

The scaled version should be $\answer{24}$ yards.

\begin{freeResponse}
\end{freeResponse}

\begin{solution}
Our scale information is that 1 ft original $=$ 4 yards scaled. We are asked to scale an
original length of 2 yards into a scaled value.

First, we need to change our 2 yards original into feet, since our scale relationship
uses feet for the ``original'' value:
\begin{gather*}
2\text{ yards original} \cdot \frac{3\text{ feet original}}{1\text{ yard original}} = 6\text{ feet original}
\end{gather*}
Now we can scale our 6 feet original length into a scaled length:
\begin{gather*}
6\text{ feet original} \cdot \frac{4\text{ yards scaled}}{1\text{ foot original}} = 24\text{ yards scaled}
\end{gather*}
So our 2 yards of the original item would be 24 yards in the scaled item.
\end{solution}
\end{problem}

Often scaling information is given as a ratio like ``this is a 1:3 scale'' which means
the scaled item is about $\tfrac{1}{3}$ smaller in length than the original. If it were
``this is a 3:1 scale'' this would mean the scaled item is about 3 times longer than
the original. Meaning the ratio is always presented as ``scale'' : ``real.''

\begin{problem}
Hot Wheels cars are generally said to be 1:64 scale compared to an actual car. If a Hot
Wheels car is 1.75 inches long, how long would you expect the real car to be? Give your
answer in feet.

The real car would be about $\answer[tolerance=0.1]{9.33}$ feet long.

\begin{solution}
1:64 scale means the real car is about 64 times longer than the Hot Wheels car, or the
Hot Wheels car is about $\tfrac{1}{64}$th the length of a real car. So if a Hot Wheels
car is 1.75 inches, we would expect the real car to be 64 times longer:
\begin{gather*}
1.75\text{ inches} \cdot 64 = 112\text{ inches} \\
112\text{ inches} \cdot \frac{1\text{ ft}}{12\text{ inches}} = \frac{112}{12}\text{ feet} = \text{about } 9.33\text{ feet}
\end{gather*}
The real car would be about 9.33 feet long.
\end{solution}
\end{problem}

\begin{problem}
If the real version of a car has 17-inch tires, a Hot Wheels version would have tires of
what diameter?

About $\answer[tolerance=0.05]{0.27}$ inches.

\begin{solution}
With a 1:64 scale, we would expect the Hot Wheels' tires to be $\tfrac{1}{64}$th the
size of the real tires. So
\begin{gather*}
17\text{ inches} \cdot \frac{1}{64} = 0.265625\text{ inches}
\end{gather*}
So we would expect the Hot Wheels' tire to be about 0.25 inches (a quarter of an inch)
in diameter.
\end{solution}
\end{problem}

\begin{problem}
Let's see what happens to the area when scaling. Suppose the Hot Wheels car is 1 in
wide --- keeping the same length of 1.75 in from above --- what area of the ground does
it cover? Give your answer in square inches.

$\answer[tolerance=0.05]{1.75}$ square inches.

\begin{solution}
The area would be the length times the width:
\begin{gather*}
1\text{ in} \cdot 1.75\text{ in} = 1.75\text{ sq in}
\end{gather*}
\end{solution}
\end{problem}

\begin{problem}
What area would the real car cover? Give your answer in square INCHES (not square
feet).

$\answer[tolerance=50]{7168}$ square inches.

\begin{solution}
We know the real car would be 112 inches long from our earlier work. Now we must find
how wide it would be by scaling the 1 in width to real length in inches, then multiply
those inch values to find the area the real car would cover:
\begin{gather*}
1\text{ inch} \cdot 64 = 64\text{ inches} \\
112\text{ inches} \cdot 64\text{ inches} = 7{,}168\text{ square inches}
\end{gather*}
\end{solution}
\end{problem}

\begin{problem}
How many times bigger is the real car's area compared to the Hot Wheels' area? How does
this compare to the length scale (1:64)?

$\answer[tolerance=20]{4096}$ times bigger.

\begin{freeResponse}
\end{freeResponse}

\begin{solution}
The Hot Wheels' area is 1.75 sq in. The real car's area is 7,168 sq in. How many times
bigger is 7,168 than 1.75? Meaning $1.75 \cdot \text{what} = 7{,}168$, so our ``what''
is $\dfrac{7{,}168}{1.75} = 4{,}096$ times bigger.

How is this value related to the 1:64 scale? $64 \cdot 64 = 4{,}096$!

\textbf{Note:} This is just like the dimension relationships, where 1 ft $=$ 12 inches
and 1 sq ft $= 12^{2} = 144$ sq inches. Our length scale factor is 1:64, but our area
scale factor is $1^{2} : 64^{2} = 1 : 4{,}096$.
\end{solution}
\end{problem}

\section*{Scaling \& Dimensions Summary}

\begin{problem}
Let's summarize what we learned in the previous problems. If something has a scale of
1:7, then

\begin{itemize}
\item 1 unit in scale $= \answer{7}$ units in real
\item 1 square unit in scale $= \answer{49}$ square units in real
\item 1 cubic unit in scale $= \answer{343}$ cubic units in real
\end{itemize}

\begin{solution}
The length factor is 7. Area picks up the square of the length factor, $7^{2} = 49$,
and volume picks up the cube of it, $7^{3} = 343$.
\end{solution}
\end{problem}

Now continue to use this pattern to help you in answers to the following problems.

\begin{problem}
A DNA helix is approximately 2 nanometers wide. We want to create a scaled up model of
DNA with a width of 5 cm. What scale would our model be (i.e.\ 100:1? 200:1?)

\hspace{1em} 1 meter $=$ 100 centimeters \qquad 1 meter $= 10^{9}$ nanometers (1 billion)

The scale would be $\answer{25000000}$ : 1.

\begin{freeResponse}
\end{freeResponse}

\begin{solution}
5 cm is how many times bigger than 2 nanometers? First we must convert the 5 cm to
nanometers:
\begin{gather*}
5\text{ cm} \cdot \frac{1\text{ meter}}{100\text{ cm}} \cdot \frac{1{,}000{,}000{,}000\text{ nm}}{1\text{ meter}}
= 50{,}000{,}000\text{ nanometers}
\end{gather*}
So 50 million nanometers is how many times bigger than 2 nanometers?
\begin{gather*}
\frac{50{,}000{,}000\text{ nm}}{2\text{ nm}} = 25{,}000{,}000\text{ times bigger}
\end{gather*}
So our scale (which is always ``scale'' : ``real'') would be 25,000,000 : 1.
\end{solution}
\end{problem}

\begin{problem}
If we want our model to be 0.5 meter tall, how many nanometers of actual DNA will be
represented in our model?

$\answer{20}$ nanometers.

\begin{solution}
We can convert the 0.5 meters to nanometers, and then reduce it by a factor of 25
million to find the ``real'' amount of nanometers:
\begin{gather*}
0.5\text{ meters} \cdot \frac{1{,}000{,}000{,}000\text{ nm}}{1\text{ meter}} \cdot \frac{1}{25{,}000{,}000}
= \frac{500{,}000{,}000}{25{,}000{,}000}\text{ nm} = 20\text{ nm}
\end{gather*}
So the 0.5 meter model would show 20 nanometers of actual DNA.
\end{solution}
\end{problem}

\begin{problem}
You've been commissioned to create a statue for a local park. You must first make a
model version to propose before the real version can be built. Suppose your statue
model is 18 inches tall and the real statue should be 15 ft tall. What is the scale of
the model as a ratio (i.e.\ 1:10? 1:25?)

1 : $\answer{10}$

\begin{solution}
How many times bigger is 15 ft than 18 inches? If we convert the 15 ft to inches and
then divide by 18 inches, we will find how many times bigger it is:
\begin{gather*}
15\text{ ft} \cdot \frac{12\text{ in}}{1\text{ ft}} = 180\text{ inches} \\
\frac{180\text{ inches}}{18\text{ inches}} = 10
\end{gather*}
So the statue will be 10 times bigger, meaning the scale (which is ``scale'' : ``real'')
would be 1:10.
\end{solution}
\end{problem}

\begin{problem}
The podium for your model is 25 square inches. What should the podium for the real
version of the statue be? Give your answer in square feet.

About $\answer[tolerance=0.5]{17.36}$ square feet.

\begin{solution}
From noticing in the Hot Wheels example that the area scale was the square of the
length scale, we can deduce that the area scale here would be $1^{2} : 10^{2} = 1 :
100$. So if the model podium is 25 sq in, the real podium will be
\begin{gather*}
25\text{ sq in} \cdot 100 = 2{,}500\text{ sq inches}.
\end{gather*}
Now we convert this to sq ft using the conversion factor 1 sq ft $= 12^{2} = 144$ sq
inches:
\begin{gather*}
2{,}500\text{ sq in} \cdot \frac{1\text{ sq ft}}{144\text{ sq in}} = \text{about } 17.36\text{ sq ft}
\end{gather*}
The podium of the real statue would be about 17.36 sq ft.
\end{solution}
\end{problem}

\begin{problem}
Suppose you used about 22 lbs of modeling clay to create your model. You find that 1 lb
of modeling clay is equal to about 16 cubic inches of clay. How much material would you
need for the real statue?

$\answer[tolerance=1000]{352000}$ cubic inches.

\begin{freeResponse}
\end{freeResponse}

\begin{solution}
Since the length scale is 1:10, the volume scale will be $1^{3} : 10^{3} = 1 : 1{,}000$.
Meaning if we know how many cubic inches are in the model, we can find out how many
cubic units should be in the real statue.

There were 22 lbs of clay used, and 1 lb $=$ 16 cu in. That means
\begin{gather*}
22\text{ lbs} \cdot \frac{16\text{ cu in}}{1\text{ lb}} = 352\text{ cu in}
\end{gather*}
We need to scale this up by 1,000, so we will need
\begin{gather*}
352\text{ cu in} \cdot 1{,}000 = 352{,}000\text{ cu inches}
\end{gather*}
So we will need 352,000 cubic inches of material for the real statue.
\end{solution}
\end{problem}

\section*{Non-Standard Unit Equivalences}

You probably think or use a sort of unit equivalence almost every day. You do this in
the form of \textbf{price}. Almost everything has a price, and the price can be thought
of as establishing an equivalence between the item or service and the \$ units. For
example, if gas is \$2.80 per gallon, this is really establishing the equivalence:
\$2.80 $=$ 1 gallon. We can then use this to convert the amount of gallons put in a car
into a cost. If you have a 15 gallon tank which you fill at \$2.80 per gallon, you are
converting 15 gallons into dollars using this equivalence:
\[
15\text{ gallons} \cdot \frac{\$2.80}{1\text{ gallon}} \;=\; \$42
\]

We can think of other relationships as equivalences as well. For example, in a previous
problem, we are told that 1 lb of modeling clay is about 16 cubic inches of clay.
Pounds and cubic inches are not measuring the same thing: pounds measure weight, cubic
inches measure volume. So generally, we cannot convert pounds to cubic inches. However,
when we know a specific substance, in this example, modeling clay, we can establish a
relationship between these two units.

\begin{problem}
In the U.S.\ we sell gas by the gallon, but most other countries sell and price it by
the liter. If you have a 15.4 gal tank in your car, how many liters is that?

\hspace{1em} 1 gal $=$ 3.78 L \qquad \$1 US $=$ \$1.38 CA

About $\answer[tolerance=0.5]{58.2}$ liters.

\begin{solution}
We want to convert the 15.4 gallons into liters, using the fact that 1 gal $=$ 3.78 L:
\begin{gather*}
15.4\text{ gal} \cdot \frac{3.78\text{ L}}{1\text{ gal}} = 58.212\text{ L}
\end{gather*}
Hence, the tank holds about 58.2 L of gas.
\end{solution}
\end{problem}

\begin{problem}
The average price per gallon is around \$3.16. How much would it cost to fill up the
15.4 gal tank?

About \$$\answer[tolerance=0.5]{48.66}$.

\begin{solution}
We want to ``convert'' the 15.4 gallons into dollars, using the fact that 1 gal $=$
\$3.16:
\begin{gather*}
15.4\text{ gal} \cdot \frac{\$3.16}{1\text{ gal}} = \$48.664
\end{gather*}
Hence, filling the tank will cost about \$48.66.
\end{solution}
\end{problem}

\begin{problem}
The average price for gas in Canada is around \$1.55 CA per liter. How much would it
cost to fill up the 15.4 gal tank?

About \$$\answer[tolerance=1]{90.21}$ CA.

\begin{solution}
We want to convert the 15.4 gallons, but measured in liters, into CA\$, using the fact
that 1 L $=$ \$1.55 CA. We found that 15.4 gal $=$ 58.2 L, so we will convert that into
\$CA:
\begin{gather*}
58.2\text{ L} \cdot \frac{\$1.55\text{ CA}}{1\text{ L}} = \$90.21\text{ CA}
\end{gather*}
Hence, filling the tank would cost about \$90.21 CA.
\end{solution}
\end{problem}

\begin{problem}
Convert the \$1.55 CA per liter into \$US per gallon. Then briefly explain why it is
important to convert both values (\$CA to \$US and liters to gallons) to compare gas
prices between here (the US) and Canada.

About \$$\answer[tolerance=0.15]{4.25}$ US per gallon.

\begin{freeResponse}
\end{freeResponse}

\begin{solution}
There are two ways to think about this.

\textbf{Method 1: convert the \$ amounts and the gal/L amounts separately.}
$\$1.55$ CA per liter is $\dfrac{\$1.55\text{ CA}}{1\text{ L}}$, so we can convert
\$1.55 CA into \$US and convert 1 L into gallons, then divide those values.
\begin{gather*}
\$1.55\text{ CA} \cdot \frac{\$1\text{ US}}{\$1.38\text{ CA}} = \$1.1232\ldots\text{ US} \\
1\text{ L} \cdot \frac{1\text{ gal}}{3.78\text{ L}} = 0.26455\text{ gal} \\
\frac{\$1.12\text{ US}}{0.265\text{ gal}} = \text{about } \$4.23\text{ per gallon}
\end{gather*}

\textbf{Method 2: convert both at the same time.} Since we know \$1 US $=$ \$1.38 CA and
1 gal $=$ 3.78 L, we can convert the units in the \$1.55 CA per liter:
\begin{gather*}
\frac{\$1.55\text{ CA}}{1\text{ L}} \cdot \frac{\$1\text{ US}}{\$1.38\text{ CA}} \cdot \frac{3.78\text{ L}}{1\text{ gal}}
= \frac{\$(1.55 \cdot 3.78)\text{ US}}{1.38\text{ gal}} = \$4.25\text{ US per gal}
\end{gather*}

Both values have to be converted because a price is a ratio of two units. Changing only
the currency, or only the volume, still leaves you comparing two quantities measured
differently --- the comparison is only meaningful once both the numerator and the
denominator are in the same units.
\end{solution}
\end{problem}

\begin{problem}
The Nisqually Glacier is a glacier on Mount Rainier in Mount Rainier National Park.
According to the National Park Service website:

\begin{quote}
``In 1857, [August] Kautz wrote in his journal that the terminus [of the Nisqually
glacier] reached a narrow rock channel located underneath the present Glacier bridge on
the road to Paradise\ldots{} By the 1950s\ldots{} the glacier had retreated nearly 1.2
miles (2 km) up past a curve in the valley.''
\end{quote}

Estimate the glacier's retreat per year. First determine it in miles per year, then
determine it in feet per year.

\hspace{1em} 1 mile $=$ 5,280 ft

About $\answer[tolerance=3]{63}$ feet per year.

\begin{freeResponse}
\end{freeResponse}

\begin{solution}
The glacier retreated about 1.2 mi from 1857 to the 1950s, which is about 100 years. So
to find the miles per year, we divide the 1.2 miles of retreat by the 100 years:
\begin{gather*}
\frac{1.2\text{ mi}}{100\text{ years}} = \frac{0.012\text{ mi}}{1\text{ year}}
\end{gather*}
To find out how many feet per year this would be, we convert the 0.012 miles to feet:
\begin{gather*}
\frac{0.012\text{ mi}}{1\text{ year}} \cdot \frac{5{,}280\text{ ft}}{1\text{ mi}}
= \frac{63.36\text{ ft}}{1\text{ year}} = \text{about 63 ft per year}
\end{gather*}
\end{solution}
\end{problem}

\begin{problem}
Why would we prefer to have the amount in feet per year instead of miles per year?

\begin{freeResponse}
\end{freeResponse}

\begin{solution}
Because most people do not have a good idea of how much 0.012 miles is, but we do have
an idea of how much 63 feet is. Since 0.012 miles $=$ 63 feet, they represent the same
length, but if we want someone to quickly understand how much this is, for a value like
this we more often measure it in feet, and so we have a better understanding of it in
feet units rather than mile units.
\end{solution}
\end{problem}

\begin{problem}
Now determine the glacier's retreat in inches per day.

About $\answer[tolerance=0.5]{2}$ inches per day.

\begin{solution}
To answer this, we must convert the 63 feet into inches, and the 1 year into days, and
then divide those values, so that we have inches per 1 day:
\begin{gather*}
\frac{63\text{ ft}}{1\text{ year}} \cdot \frac{12\text{ in}}{1\text{ ft}} \cdot \frac{1\text{ year}}{365\text{ days}}
= \frac{2.0712\ldots\text{ in}}{1\text{ day}} = \text{about 2 inches per day}
\end{gather*}
The glacier retreats an average of 2 inches per day.

\textbf{Note:} You could also do these conversions (ft to inches, and years to days)
separately, like we did in Method 1 of the gas price problem.
\end{solution}
\end{problem}

\begin{problem}
Which is a ``better'' way of describing the glacier's retreat: feet per year or inches
per day? Support your choice based on math work and based on how a glacier might
actually retreat in the ``real world.''

\begin{freeResponse}
\end{freeResponse}

\begin{solution}
Based on the math, both descriptions are about the same. You could say the 63 ft/year is
a bit more accurate, as we did some more rounding to get the 2 inches per day and more
rounding gets us further from the ``true'' value. However, mathematically, these are
essentially the same.

Based on how a glacier might actually retreat in the ``real world,'' the 63 ft per year
would probably be preferable over the 2 inches per day value. If you actually stood at
the edge of the glacier for a day, you would not see 2 inches of the glacier melt away.
In the summer, the glacier may be melting more than 2 inches per day, and in the winter
the glacier might actually gain some length from more snow and freezing temperatures,
but from the start to the finish of a year you would see it receding. Even from year to
year there would be variation, but the 63 ft per year would give a better idea of what
is happening with the glacier.
\end{solution}
\end{problem}

\end{document}
